% =====================================================================
% DBA - Proyecto final - Entrega 3
% Huellas de edificios: descarga, limpieza, carga en MongoDB y EDA
% =====================================================================
\documentclass[11pt,a4paper]{article}
% =====================================================================
% Preámbulo compartido - DBA Proyecto final
% =====================================================================
\usepackage[utf8]{inputenc}
\usepackage[T1]{fontenc}
\IfFileExists{lmodern.sty}{\usepackage{lmodern}}{}
\usepackage[spanish,es-noshorthands,es-tabla]{babel}
\usepackage[margin=2.5cm]{geometry}
\usepackage{amsmath,amssymb}
\usepackage{booktabs,tabularx,array}
\usepackage{enumitem}
\usepackage[dvipsnames,table]{xcolor}
\usepackage{graphicx}
\usepackage{tikz}
\usetikzlibrary{calc,positioning}
\usepackage{listings}
\usepackage{caption}
\usepackage{float}
\usepackage[hidelinks]{hyperref}
\usepackage{xurl}

% --- Colores --------------------------------------------------------
\definecolor{acento}{HTML}{2F6FDB}
\definecolor{gris}{HTML}{6B6B6B}
\hypersetup{colorlinks=true, linkcolor=acento, urlcolor=acento, citecolor=acento}

% --- Tablas ---------------------------------------------------------
\newcolumntype{L}{>{\raggedright\arraybackslash}X}
\renewcommand{\arraystretch}{1.25}

% --- Listas ---------------------------------------------------------
\setlist{itemsep=2pt, topsep=4pt}

% --- Código (JSON / mongosh) ----------------------------------------
\lstdefinestyle{codigo}{
  basicstyle=\ttfamily\footnotesize,
  breaklines=true,
  columns=fullflexible,
  keepspaces=true,
  frame=single,
  framerule=0.3pt,
  rulecolor=\color{gray!50},
  backgroundcolor=\color{gray!6},
  xleftmargin=4pt, xrightmargin=4pt,
  showstringspaces=false,
  literate={á}{{\'a}}1 {é}{{\'e}}1 {í}{{\'i}}1 {ó}{{\'o}}1 {ú}{{\'u}}1
           {Á}{{\'A}}1 {É}{{\'E}}1 {Í}{{\'I}}1 {Ó}{{\'O}}1 {Ú}{{\'U}}1
           {ñ}{{\~n}}1 {Ñ}{{\~N}}1
}
\lstset{style=codigo}

% --- Atajos ---------------------------------------------------------
\newcommand{\campo}[1]{\texttt{#1}}
\newcommand{\pendiente}[1]{\textcolor{gris}{\textit{[Pendiente -- #1]}}}
% Ruta de archivo larga: igual que \campo pero se puede partir en / y _
% (xurl ya amplía los puntos de corte de \url para que esto funcione).
\newcommand{\ruta}[1]{\begingroup\urlstyle{tt}\url{#1}\endgroup}

\begin{document}

% ------------------------------- Portada -----------------------------
\begin{titlepage}
  \centering
  {\large Pontificia Universidad Javeriana\par}
  {\large Departamento de Ingeniería de Sistemas\par}
  \vspace{0.8cm}
  {\large Administración de Bases de Datos (DBA) -- 2026-30\par}
  \vspace{2.5cm}
  {\Huge\bfseries Proyecto final\par}
  \vspace{1cm}
  {\LARGE Entrega 3: Huellas de edificios -- descarga,\\ carga en MongoDB y EDA\par}
  \vspace{1cm}
  {\large Potencial solar en techos de los municipios PDET\par}
  \vspace{3cm}
  {\large
    Maria Camila Beltran \\
    \vspace{0.5cm}
    Jorge Fortich \\
    \vspace{0.5cm}
    Jose Eduardo Valeriano \\
    \vspace{0.5cm}
    Samuel Osorio Rodríguez\par}
  \vspace{3cm}
  {\large Profesor: Andrés Oswaldo Calderón Romero, Ph.D.\par}
  \vspace{3cm}
  {\large {Noviembre 16 de 2026}\par}
\end{titlepage}

\tableofcontents
\newpage

% =====================================================================
% 1. Contexto
% =====================================================================
\section{Contexto}
\label{sec:contexto}

Esta entrega corresponde a la \textbf{semana 3} del plan de implementación
definido en la Entrega 1: descargar las huellas de edificios de Google Open
Buildings v3 y Microsoft Building Footprints, quedarse con las que caen en los
170 municipios PDET, limpiarlas y cargarlas en MongoDB con índices espaciales.
La asignación de cada edificio a su municipio y el cálculo de área por
municipio son de la semana 4.

El flujo de cada fuente es el mismo y lo encadena \campo{src/edificios/run.py}
(lo que corre \campo{make edificios}): descarga con \campo{SHA-256}, prefiltro
a PDET por la regla del centroide, limpieza compartida
(\campo{src/edificios/comun.py}) hacia un Parquet con el mismo contrato de
columnas, y finalmente creación de colecciones, carga e índices en MongoDB.

\begin{table}[H]
  \centering
  \caption{Cifras principales de la entrega.}
  \label{tab:resumen}
  \begin{tabularx}{\textwidth}{L r r}
    \toprule
    & Google & Microsoft \\
    \midrule
    Archivos descargados & 8 celdas S2 (3,11\,GB) & 125 cuadros quadkey (373,4\,MB) \\
    Edificios leídos del área descargada & 16\,512\,141 (dentro de Colombia) & 4\,595\,615 \\
    Edificios en PDET (en la colección) & 2\,694\,786 & 1\,761\,048 \\
    Polígonos inválidos reparados (GEOS) & 0 & 28 \\
    Duplicados por \campo{geom\_hash} & 0 & 0 \\
    \bottomrule
  \end{tabularx}
\end{table}

% =====================================================================
% 2. Plan de implementación
% =====================================================================
\section{Plan de implementación}
\label{sec:plan}

\begin{table}[H]
  \centering
  \caption{Plan de implementación y estado.}
  \label{tab:plan}
  \begin{tabularx}{\textwidth}{l L L l}
    \toprule
    Semana & Fase & Resultado principal & Estado \\
    \midrule
    1 & Infraestructura & MongoDB 8 en Docker, repositorio y configuración & Hecha \\
    2 & Límites PDET & 170 municipios cargados, validados e indexados & Hecha \\
    3 & Huellas de edificios & Google y Microsoft descargados, limpios, cargados e
        indexados; EDA breve & \textbf{Esta entrega} \\
    4 & Análisis e informe & Asignación por centroide, agregación en
        \campo{resultados}, mapas y ranking & Pendiente \\
    \bottomrule
  \end{tabularx}
\end{table}

\textbf{Organización del equipo.} Una rama por persona y por
tarea, con revisión cruzada antes del merge a \campo{main}. Reparto de esta
semana: descarga y limpieza de Google (Jorge); descarga y limpieza de
Microsoft (Maria Camila); esquema, carga, índices y \campo{make edificios}
(Samuel); EDA (Jose Eduardo).

% =====================================================================
% 3. Modelo de datos
% =====================================================================
\section{Modelo de datos}
\label{sec:modelo}

Las colecciones \campo{edificios\_google} y \campo{edificios\_microsoft}
comparten esquema (Entrega 1, sección 5). Entre las dos entregas se hizo un
ajuste: el Parquet intermedio no lleva \campo{cod\_mpio} ni \campo{lote\_id};
el segundo lo pone la carga y el primero queda en \campo{null} hasta la
semana 4.

\pendiente{Tabla del diccionario de \campo{edificios\_*} y diagrama E-R
(se puede reutilizar \ruta{../Entrega 2/img/diagrama_er.png}); Samuel, sección
\ref{sec:samuel3-schema}, ya describe el validador.}

% =====================================================================
% 4. Google (Jorge)
% =====================================================================
\section{Google Open Buildings: descarga y limpieza}
\label{sec:google}

\pendiente{Jorge. Material disponible: \campo{descarga\_google.py},
\campo{limpieza\_google.py}, \campo{control\_google.tex}. Cifras ya verificadas
en dos máquinas: 8 celdas S2 con \texttt{--solo-pdet} (3,11\,GB, SHA-256
\texttt{19e8f1e7\ldots}), 31\,370\,317 en el rectángulo de Colombia,
16\,512\,141 dentro de Colombia, 2\,694\,786 tras el prefiltro PDET, 0
inválidos, 0 descartados, 0 duplicados. Confianza 0,65--0,9834.}

% =====================================================================
% 5. Microsoft (Maria Camila)
% =====================================================================
\section{Microsoft Building Footprints: descarga y limpieza}
\label{sec:microsoft}

\pendiente{Maria Camila. Material disponible: \campo{descarga\_microsoft.py},
\campo{limpieza\_microsoft.py}, \campo{control\_microsoft.csv}. Cifras
verificadas: 232 cuadros de Colombia en el índice, 125 tocan PDET (373,4\,MB,
SHA-256 global \texttt{3ca8b1d5\ldots}), 4\,595\,615 leídos, 2\,834\,567 fuera de
PDET, 1\,761\,048 dentro, 28 inválidos reparados, 0 duplicados.}

% =====================================================================
% 6. Esquema, carga e índices (Samuel)
% =====================================================================
\newcommand{\figsamuel}[2]{%
  \IfFileExists{img/#1.png}%
    {\includegraphics[width=\textwidth]{img/#1.png}}%
    {\pendiente{captura \texttt{#2}}}}

\section{Esquema, carga e índices de edificios en MongoDB}
\label{sec:samuel3}

Esta sección cubre \campo{schema.py}, \campo{carga.py}, \campo{indices.py},
\campo{evidencia.py} y \campo{run.py}, todos en \campo{src/edificios/}. Toman
los dos Parquet limpios (uno por fuente, con el mismo contrato de columnas) y
los dejan en MongoDB listos para la asignación por centroide de la semana~4.

\subsection{Colecciones con validación \texorpdfstring{\campo{\$jsonSchema}}{\$jsonSchema}}
\label{sec:samuel3-schema}

\campo{schema.py} crea \campo{edificios\_google} y \campo{edificios\_microsoft}
con el mismo validador, derivado del diccionario de datos de la Entrega~1
(sección~5). Campos requeridos: \campo{fuente\_id}, \campo{geometry}
(\campo{Polygon}), \campo{centroide} (\campo{Point} de dos \campo{double}),
\campo{area\_m2} ($\geq 0$), \campo{geom\_hash} y \campo{lote\_id}.
\campo{confianza} (entre 0 y 1), \campo{area\_m2\_fuente} y \campo{cod\_mpio}
aceptan \campo{null}: Microsoft no entrega confianza ni área propia, y
\campo{cod\_mpio} se llena en la semana~4. La validación es estricta
(\campo{validationLevel="strict"}), así que MongoDB rechaza al insertar
cualquier documento que no cumpla. Igual que en la Entrega~2, si la colección
ya existe el script no modifica su validador.

\begin{figure}[H]
  \centering
  \figsamuel{cap_schema_edificios}{schema.py}
  \caption{Ejecución de \campo{schema.py}: se crean las dos colecciones de
  edificios con validación \campo{\$jsonSchema}.}
  \label{fig:cap-schema-edif}
\end{figure}

\begin{figure}[H]
  \centering
  \figsamuel{cap_validacion_edificios}{validación}
  \caption{Un documento de edificio sin campos obligatorios es rechazado por
  el validador (código de error 121).}
  \label{fig:cap-validacion-edif}
\end{figure}

\subsection{Carga por lotes e idempotente}
\label{sec:samuel3-carga}

\campo{carga.py} lee el Parquet con \campo{pyarrow} en lotes de 5\,000 filas
(\campo{edificios.tamano\_lote\_insercion} en \campo{config.yaml}), de modo que
nunca se sube el archivo completo a memoria. Cada lote se convierte en
documentos: la geometría pasa de WKB a GeoJSON con Shapely, el centroide se
arma como \campo{Point} y \campo{cod\_mpio} queda en \campo{null}. La
inserción usa \campo{insert\_many(ordered=False)}, así un documento rechazado
no frena el resto del lote.

Decisiones de diseño:
\begin{itemize}
  \item \textbf{Idempotencia.} Antes de cargar, borra los documentos previos
        de la misma fuente; correr la carga dos veces deja el mismo conteo
        (comprobado: segunda corrida con \campo{make edificios}, mismas
        cifras).
  \item \textbf{\campo{lote\_id}.} Cada ejecución genera un \campo{ObjectId}
        que se guarda en cada edificio y es el \campo{\_id} de su fila en
        \campo{bitacora\_cargas}, con inicio, fin, registros correctos,
        rechazados y velocidad. Así cada edificio se rastrea hasta la carga
        que lo trajo.
  \item \textbf{Índices.} Se borran los índices (salvo \campo{\_id}) antes de
        cargar y los de \campo{centroide}, \campo{cod\_mpio} y
        \campo{geom\_hash} se crean después, porque mantenerlos durante la
        inserción es más lento. El de \campo{geometry} es la excepción, por
        la razón de la sección~\ref{sec:samuel3-s2}.
  \item Se actualiza \campo{fuentes.n\_registros} con el total cargado.
\end{itemize}

\begin{table}[H]
\centering\small
\begin{tabularx}{\textwidth}{@{}lrrrrL@{}}
\toprule
\textbf{Fuente} & \textbf{Documentos} & \textbf{Rechazados} &
\textbf{Reparados (S2)} & \textbf{Tiempo (s)} & \textbf{Velocidad} \\
\midrule
Google    & 2\,694\,786 & 0 & 0  & 63,5 & 42\,468 docs/s \\
Microsoft & 1\,761\,048 & 0 & 29 & 43,4 & 40\,608 docs/s \\
\bottomrule
\end{tabularx}
\caption{Carga a MongoDB (corrida completa de \campo{make edificios}, MongoDB
8 en Docker local, un solo equipo). Los tiempos dependen de la máquina.}
\label{tab:carga-edif}
\end{table}

\begin{figure}[H]
  \centering
  \figsamuel{cap_carga_edificios}{carga.py}
  \caption{Ejecución de \campo{carga.py} para las dos fuentes.}
  \label{fig:cap-carga-edif}
\end{figure}

\subsection{Un hallazgo: polígonos válidos para GEOS que MongoDB rechaza}
\label{sec:samuel3-s2}

La primera versión creaba todos los índices después de la carga. El índice
\campo{2dsphere} sobre \campo{geometry} de Microsoft falló: MongoDB construye
sus claves con geometría esférica (biblioteca S2) y encontró un polígono cuyas
aristas se cruzan (\emph{``Loop is not valid\ldots Edges 1 and 3 cross''}),
aunque la limpieza ya lo había dado por válido con GEOS, que trabaja en el
plano. Con una sola geometría así, la construcción de todo el índice aborta.

Se midió con una muestra de 300\,000 edificios: 5 rechazados, es decir unos
2 por millón. La solución fue crear el índice de \campo{geometry} \emph{antes}
de cargar: con el índice activo MongoDB rechaza solo los documentos afectados
y la carga sigue (en la muestra costó 5,5 s). Cada rechazado se repara
redondeando las coordenadas a $10^{-6}$ grados (unos 11\,cm) con
\campo{shapely.set\_precision} y se reintenta; como respaldo existe el casco
convexo. En una prueba con los 5 casos de la muestra, el redondeo los resolvió
todos, mientras que \campo{buffer(0)} y \campo{simplify} no resolvieron
ninguno. En la carga completa, Microsoft tuvo 29 reparados y Google ninguno,
y no se descartó ningún edificio. Los reparados conservan el
\campo{geom\_hash} del polígono original, y el conteo de reparados queda en
\campo{bitacora\_cargas.parametros.reparados\_s2}.

\subsection{Índices}
\label{sec:samuel3-indices}

\campo{indices.py} crea los índices del diccionario de datos (Entrega~1,
sección~5.5) en las dos colecciones y mide cuánto tarda cada uno:

\begin{table}[H]
\centering\small
\begin{tabularx}{\textwidth}{@{}p{3.6cm}Lrr@{}}
\toprule
\textbf{Índice} & \textbf{Para qué} & \textbf{Google (s)} & \textbf{Microsoft (s)} \\
\midrule
\campo{2dsphere} en \campo{centroide} & \campo{\$geoWithin} de la asignación (semana~4) & 7,0 & 4,4 \\
\campo{2dsphere} en \campo{geometry} & \campo{\$geoIntersects} en límites & \multicolumn{2}{c}{creado antes de cargar} \\
\campo{cod\_mpio} & agregación por municipio & 4,7 & 1,4 \\
\campo{geom\_hash} (único) & evitar duplicados & 3,5 & 2,3 \\
\bottomrule
\end{tabularx}
\caption{Tiempo de creación de cada índice.}
\label{tab:indices-edif}
\end{table}

El índice único sobre \campo{geom\_hash} es una segunda barrera contra
duplicados, además del \campo{DISTINCT ON (geom\_hash)} de la limpieza: si
algún duplicado llegara a la colección, la creación del índice fallaría.

\begin{figure}[H]
  \centering
  \figsamuel{cap_indices_edificios}{indices.py}
  \caption{Ejecución de \campo{indices.py}.}
  \label{fig:cap-indices-edif}
\end{figure}

\subsection{Evidencia de que los índices se usan}
\label{sec:samuel3-evidencia}

\campo{evidencia.py} usa como ventana el polígono de un municipio PDET
(Caldono, 19137) y corre dos consultas por fuente. Con \campo{.explain()}
exige que el plan use \campo{IXSCAN} y no \campo{COLLSCAN}, y para
contrastar repite la consulta de la semana~4 forzando el recorrido completo
con \campo{hint(\{\$natural: 1\})}.

\begin{table}[H]
\centering\small
\begin{tabularx}{\textwidth}{@{}llLrr@{}}
\toprule
\textbf{Fuente} & \textbf{Plan} & \textbf{Consulta} & \textbf{Resultados} & \textbf{Tiempo (s)} \\
\midrule
Google    & IXSCAN   & \campo{\$geoWithin} sobre \campo{centroide}   & 15\,303 & 0,44 \\
Google    & IXSCAN   & \campo{\$geoIntersects} sobre \campo{geometry} & 15\,304 & 0,37 \\
Google    & COLLSCAN & \campo{\$geoWithin} sin índice                 & 15\,303 & 2,97 \\
Microsoft & IXSCAN   & \campo{\$geoWithin} sobre \campo{centroide}   & 12\,547 & 0,33 \\
Microsoft & IXSCAN   & \campo{\$geoIntersects} sobre \campo{geometry} & 12\,548 & 0,29 \\
Microsoft & COLLSCAN & \campo{\$geoWithin} sin índice                 & 12\,547 & 1,95 \\
\bottomrule
\end{tabularx}
\caption{Planes de ejecución y tiempos en Caldono. Es una sola consulta por
caso, así que los tiempos son indicativos, no un \emph{benchmark}.}
\label{tab:evidencia-edif}
\end{table}

Con el índice, la consulta de asignación corre unas 6,7 veces más rápido en
Google y unas 5,9 en Microsoft que sin él, con el mismo resultado. La
diferencia de un edificio entre \campo{\$geoWithin} (centroide dentro) y
\campo{\$geoIntersects} (huella que toca el municipio) corresponde a un
edificio que cruza el límite: es la razón de usar el centroide como regla de
asignación, para que cada edificio pertenezca a un solo municipio.

\begin{table}[H]
\centering\small
\begin{tabularx}{\textwidth}{@{}lrrrL@{}}
\toprule
\textbf{Colección} & \textbf{Documentos} & \textbf{Datos (MB)} & \textbf{Índices (MB)} & \textbf{Mayor índice} \\
\midrule
\campo{edificios\_google}    & 2\,694\,786 & 1\,279 & 251 & \campo{geom\_hash} (111 MB) \\
\campo{edificios\_microsoft} & 1\,761\,048 & 835    & 194 & \campo{geom\_hash} (88 MB) \\
\bottomrule
\end{tabularx}
\caption{Tamaño de las colecciones (\campo{collStats}).}
\label{tab:tamano-edif}
\end{table}

\begin{figure}[H]
  \centering
  \figsamuel{cap_evidencia_edificios}{evidencia.py}
  \caption{Salida de \campo{evidencia.py}: planes con \campo{IXSCAN} y tamaños.}
  \label{fig:cap-evidencia-edif}
\end{figure}

\subsection{Ejecución de punta a punta: \texorpdfstring{\campo{make edificios}}{make edificios}}
\label{sec:samuel3-run}

\campo{run.py} encadena, para cada fuente, descarga, limpieza, creación de
colecciones, carga, índices y evidencia, con un solo comando. Cada paso es
reanudable: la descarga reutiliza los archivos ya bajados y recalcula el
SHA-256 de la fuente; la limpieza regenera el Parquet y la carga es idempotente. Se corrió
\campo{make edificios} sobre datos ya descargados y se obtuvieron
exactamente las mismas cifras que en la primera ejecución (2\,694\,786 y
1\,761\,048 documentos, 29 reparados en Microsoft). Esto verifica el flujo y
la idempotencia; \textbf{no} es una descarga desde cero, que se probó
por separado.\footnote{\pendiente{confirmar tras borrar \texttt{data/raw/} y correr
\texttt{make edificios}; ajustar o quitar esta nota.}}

\begin{figure}[H]
  \centering
  \figsamuel{cap_run_edificios}{make edificios}
  \caption{Final de la ejecución de \campo{make edificios}.}
  \label{fig:cap-run-edif}
\end{figure}


% =====================================================================
% 7. EDA (Jose Eduardo)
% =====================================================================
\section{Análisis exploratorio breve}
\label{sec:eda}

\pendiente{Jose Eduardo. Material de contraste ya calculado: Google tiene unas
1,53 veces más edificios que Microsoft pero más pequeños (media 82,4 frente a
125,6\,m$^2$; mediana 57,9 frente a 77,0), y el área total es parecida
($\approx$222 frente a $\approx$221\,km$^2$, calcular la suma exacta).
Diferencias menores entre pipelines a mencionar: Google trae una columna
\campo{cod\_mpio} vacía; hash MD5 frente a SHA-1; orientación de anillos;
cálculo del centroide.}

% =====================================================================
% 8. Conclusiones
% =====================================================================
\section{Conclusiones y próximos pasos}
\label{sec:conclusiones}

\pendiente{Redactar al final, con las cifras definitivas. Puntos a cubrir:
las dos fuentes quedan cargadas y reproducibles con \campo{make edificios};
el hallazgo de la validación esférica de MongoDB
(sección~\ref{sec:samuel3-s2}); y el paso a la semana 4: asignar cada edificio
a su municipio con \campo{\$geoWithin} sobre \campo{centroide}, agregar en
\campo{resultados} y producir las tablas y mapas finales.}

\end{document}
